% 原书第 107 页（PDF 第 120 页）
\BookPage{107}{120}
\BookHeading{清代扬州的盐务官述略}{清代扬州的盐务官述略}{杜家骥*}
扬州因盐业的发展而繁盛。清代，两淮盐区是全国最大的食盐产、销区，盐引销售范围达及江苏、安徽、江西、湖北、湖南及河南等省。盐课税收占全国盐税收入的一半左右。而管理两淮盐区产、运、销以保障其盐课收入的两个专职官一盐政、盐运使，及其衙署的其他盐务官吏，便设在扬州。盐政衙署在开明桥东南（扬州市新华中学内），盐运使衙署在大东门外（扬州市人民政府所在地），史称「广陵鹾署」2。

清代沿袭历代王朝制度，盐业由国家垄断，设官管理。管理形式有几种，而实行地域范围最广、时间最长的管理形式是官督商销。官方对盐的控制性专卖，主要是引盐制。盐引，是国家发给盐商收盐、售盐的凭证，盐商领取官引，按引向国家纳课。每引有额定斤数，引外超出斤数之盐、用残引（用过之引）行销之盐，以及无引之盐，均为私盐，销售私盐，国家以法律惩处。盐商有散商，有承包较多盐引而分予各盐商的「总商」。各盐区之盐有划定的销售区，称为「引岸」，盐商须在规定引岸内销售。官方管理盐的产、运、销，掌握盐价的制定，缉拏私盐之行销，构成官督商销的主要内容。

管理盐务的官员众多，既有专职盐务官，又有兼管官。兼管官有各销盐省的总督或巡抚，其属下地方官也需配合，其中最重要的官员是盐法道（或称道员）。盐法道，是按省设，多为兼职，有兼其他专职的盐茶道、盐驿道、盐粮道，有兼地方

\vfill
\begin{BookNotes}
\BookNote{*}{作者任职于南开大学中国社会史研究中心。}
\BookNote{1.}{嘉庆《两淮盐法志》，卷55〈杂记·碑刻〉，李发元〈盐院题名碑记〉：「两淮岁课，当天下租庸之半。」并见陈锋，《清代盐政与盐税》第171页的表42「两淮与全国税课岁入得比较」。中州古籍出版社，1988年版。}
\BookNote{2.}{见《扬州古港史》第五章第四节之五「扬州港盐运的管理」。《扬州古港史》编写组吴家兴等集体编写，人民交通出版社，1988年版。}
\end{BookNotes}

% 原书第 108 页（PDF 第 121 页）
\BookPage{108}{121}
\noindent 区划道事务者，还有兼盐运使的。销盐省的盐法道，主要掌引盐的疏通销售、缉私之事。

有关两淮盐务的研究，成果较多，如（日）佐伯富《清代盐政之研究》3、徐泓《清代两淮盐场的研究》4为这方面的重要成果。两淮盐区盐务官的研究，仅据所查到的介绍如下。张荣生〈古代淮南盐区的盐官制度〉一文，在简要介绍淮南盐官的设置、职掌、俸给之外，详细论述了盐课大使这一职官职权的扩大，以及盐官的选任、考核。吴静芳〈清代前期巡视两淮盐政官员之探讨〉6，考察了清代前期巡视两淮盐政官员的任命、其职衔的变迁、内务府包衣之充任两淮盐政，以及两淮盐政与户部、地方官、盐商的关系。何峰〈明清淮南盐区盐场大使的设置、职责及其与州县官的关系〉一文，对淮南盐区盐场大使的职掌、与州县官的关系，有更深入的考察与论述。曹爱生〈清代两淮盐官制度〉，着重从税收的角度论述对两淮盐官的选拔使用及监督、考核问题。此外，前引佐伯富〈清代盐政之研究〉、徐泓《清代两淮盐场的研究》、陈锋《清代盐政与盐税》，及李克毅〈清代盐官与盐政〉，也涉及两淮盐官的问题。本人于两淮盐官素无研究，今参加「盛清社会与扬州」学术会议，仅将设于扬州的两淮盐务官与其下属盐官的系统设置，以及盐政等重要盐官的职掌、有关活动等作简要介绍，希望能对前人的考察有所补充。

一、两淮盐区专职盐务官的设置与其职掌

两淮盐政。盐政按盐区设置，是盐区最高专职盐务官，总掌该盐区的产、运、销及相关事务。最初主要由监察御史充任，所以称「巡盐御史」，后来选用范围扩大，尤其是内务府官员充任者较多，也挂监察御史衔。又因监察御史是都察院官，

\vfill
\begin{BookNotes}
\BookNote{3.}{佐伯富，〈清代盐政之研究〉（昭和31年）。顾南、顾学稼译文，连载于《盐业史研究》1993年第2期以后。}
\BookNote{4.}{徐泓，《清代两淮盐场的研究》，（台北）嘉新水泥公司文化基金会，1972年。}
\BookNote{5.}{文载《盐业史研究》2001年3期。}
\BookNote{6.}{东海大学历史研究所硕士论文，2004年。}
\BookNote{7.}{文载《盐业史研究》2006年1期。}
\BookNote{8.}{文载《盐业史研究》2006年2期。}
\BookNote{9.}{文载《中国社会经济史研究》1990年4期。}
\end{BookNotes}

% 原书第 109 页（PDF 第 122 页）
\BookPage{109}{122}
\noindent 所以巡盐御史及后来所称的「盐政」，又有「盐院」之称。从现在所掌握的资料看，至晚在雍正初年，已有巡盐御史称盐政的情况，大致在乾隆前期以后，巡盐御史都已习称「盐政」10。巡盐御史一盐政属于「钦差」性质的官员，为「一年期满之差」，由都察院将该院的监察御史、给事中名单开列给皇帝，由皇帝在其中圈定某人充任，皇帝也可以特旨的形式，差任内务府官员等担任11。盐政因是差官，担任者品级不同，所以没有专定品级，但因为属于「钦差」，为盐区最高专职盐务官，虽然充任者多三品以下，五品官居多，但地位高，本职权力较大。

两淮盐政的具体职掌是：巡视两淮盐课，督销江苏、安徽、江西、湖南、湖

\vfill
\begin{BookNotes}
\BookNote{10.}{「盐政一词有两个含义，一是指盐法政务，也可简称为盐务；一是作为巡盐御史的称谓，为官称。后者出现较晚，但究竟何时始出现，进一步说，作为官称的盐政最早何时出现，目前尚无专门研究，说法也不一。有的说是康熙中（见前引〈清代的盐官与盐政〉一文）或康熙末（见前引〈清代前期巡视两淮盐政官员之探讨〉一文），有的认为雍正年间改称盐政（见前引〈古代淮南盐区的盐官制度〉一文），有的说乾隆后简称盐政（见前引《清代盐政与盐税》一书）。盖因两种含义的「盐政均为名词，在史料中出现容易造成理解上的差异，如雍正《大清会典》卷51〈户部二十九〉页18（殿本）记：「顺治十三年题准，盐政衙门不许商役互充」，其中「盐政衙门」一词中的「盐政」，就容易误解，我认为应指的是盐务，而不是官称。前述诸说，只有清代前期巡视两淮盐政官员之探讨〉一文曾举史料作为根据，认为康熙末开始出现盐政头衔，举出光绪《大清会典事例》卷230，页9-10以下记载为据：康熙四十四年题准，盐政文移，向由都察院转行转送。嗣后盐政一应事宜，有关各部者，该部径自行文该差，该差亦自行咨呈达部。仍具呈都察院存验。(此段资料也见于雍正《大清会典》卷51〈户部二十九〉页15（殿本)）。这段史料颇费解，或者说是容易误解。主要是「盐政」一词，若单从文法上讲，理解为是官称也通。但若结合制度，就不大像是指的官称，理解为是「盐务」是否更符合原意，因为下文又专门提到巡盐御史为「该差」，把他作为差职的官称。如果先已称巡盐御史为盐政了，下文就仍称其为盐政，而不称其为「该差」了。这里的「该差」，是指巡视盐务之差，也即作为差职的巡盐御史，后文仍是延续前文之意，继述盐务及盐务文移这一事务。这段资料的意思，似应是「盐务文移，向由都察院转行转送。嗣后盐务一应事宜，有关各部者，该部径自行文该盐务的差官巡盐御史，该差官巡盐御史亦自行咨呈达部。仍具呈都察院存验」。之所以作这样的判断，还因迄未发现康熙后期巡盐御史称盐政的情况，包括盐官自己在奏章中的职称。现在所见最早的官称盐政，是雍正初年，如《清世宗实录》卷16，雍正二年二月清世宗训谕各省关差、盐差等官的训词中有「盐政不得罔恤穷商，独专厚利，硬派州县，计口征钱」之语。卷19还有雍正二年闰四月「升两淮盐政右佥都御史谢赐履为都察院左副都御史」之语。若以《清世宗实录》修于乾隆前期，所用官称之词不一定能说明问题的话，那么，成书于雍正十年记事止于雍正五年的雍正《大清会典》，则其卷49〈户部二十七〉页12（殿本）记有：「雍正三年覆准…查明灶户名姓，开载图籍，申详该盐政及该督抚存案。该盐政造册送部查核。」卷51〈户部二十九〉页29（殿本）有「雍正三年覆准，嗣后拿获私贩…或有不肖场官刻剥灶户，借申详为恐吓之端，该盐政即行题参的记载。至于在雍正初年以前乃至康熙朝是否有官称盐政的资料，目前因未发现，不能断定，现在只能说「至晚在雍正初年」。另外，从大量的资料来看，乾隆前期，巡盐御史渐为盐政之称取代。}
\BookNote{11.}{嘉庆《大清会典》卷6〈吏部〉，第631册226页，台北文海出版社影印本。}
\end{BookNotes}

% 原书第 110 页（PDF 第 123 页）
\BookPage{110}{123}
\noindent 北、河南各府州县额定盐引，审定盐价，察照下设之盐运使司、各分司及制盐场官丁亭户，严行卫所、有司缉捕私盐之贩卖1²。巡盐御史一盐政是皇帝钦派的督察盐务的差官，其职掌的性质特点是督察，两淮盐政图明阿在其到任不久给乾隆皇帝的奏折中，就汇报自己在督察盐官、盐商、销引等几方面的工作:「奴才到任后，查盐属各官，惟盐运使仓圣裔近在同城，连日见面，奴才看其才质谨朴、明晰盐务，尚属认真办事之人。其余监掣同知，暨各分司大使、坝官等，奴才由淮至扬，亦俱陆续接见，询以所管事宜，大约皆循分供职。奴才惟留心查察，随事督率，倘其中有于官方政体稍涉怠玩陨越，即当据实查办，恭请圣训…（两淮商众），奴才于其公同进见时，询问盐务情形及应办事件，即与谆切告诫…又查庚子纲盐自上年五月开纲至今，已办运九十三万余引，瞬届奏销，即须开运辛丑新纲，奴才督同运司攒旧办新，俾商运流通…。」13另外，是巡查产盐场及附近境地，两淮巡盐御史李煦给皇帝的奏折，便述其「公事稍瑕，即轻骑巡查三江营等处」私枭，又「历海边各场，细为巡查沿海灶户」14。两淮盐政还兼与督抚掌运判之选任，及所属盐官之考核15。此外，两淮盐政还承担某些特殊事务，如掌管行宫、皇帝南巡之接驾，及诸文化事务等，见后述。两淮盐政衙署中没有下属官，而设置非官员的书吏，称为「典吏」，额设十二人，协助盐政办理往来公文及各类册簿、文字杂务。

各盐区盐政，后来逐渐取消，改由督抚兼理，两淮盐政取消较晚，是在道光十年，由两江总督兼理。实际取消，有可能是在道光十一年初，因该动议的奏报时间是道光十年十二月十三日，皇帝朱批「另有旨」，迨圣旨发下、传至两淮，再实际落实，需要一段时间。取消的原因，该奏折称，是「总缘盐政无管辖地方之责，文武员弁既非所属，而疏销责在州县，巡缉责在弁兵，难期令行禁止。且一年即应候代，虽间有蒙恩留任数年，而所属官、吏、商、丁鲜不存任满当去之见。以故诸事废弛，浮费重叠，日甚一日。此次虽大加厘剔，若仍由盐臣经理，恐不能破除积

\vfill
\begin{BookNotes}
\BookNote{12.}{见光绪重修《两淮盐法志》卷129职官门·官制上〉，及乾隆《历代职官表》卷61〈盐政·国朝官制·盐政〉，下册第1173页，上海古籍出版社，1989年影印武英殿本。}
\BookNote{13.}{中国第一历史档案馆、扬州市档案局编《清宫扬州御档选编》第3册第196-198页，图明阿奏折。广陵书社，2009年版。}
\BookNote{14.}{《清宫扬州御档选编》第1册第32-33页，两淮巡盐御史李煦奏折。}
\BookNote{15.}{这方面，前引张荣生〈古代淮南盐区的盐官制度〉、曹爱生〈清代两淮盐官制度〉两文有较详细叙述，请参考。}
\end{BookNotes}

% 原书第 111 页（PDF 第 124 页）
\BookPage{111}{124}
\noindent 习、力振因循…所有两淮盐务，似亦应查照户部前奏，将盐政裁撤，责成地方官办理，则事权归一，呼应更灵，即年例开销各项浮费必可大加裁省，实与盐务、地方均有裨益。相应请旨，将两淮盐务改归两江总督管理，其盐政一缺拟请裁撤…道光十年十二月十三日」16。

盐运使，从三品，每盐区设一人。两淮盐运使掌督察两淮盐场之制盐，稽核所派销之盐引，兼理所管行盐州县，疏销盐引，以保证盐课之征收。带管盐义仓及下河水利相关之事7。道光年间撤销盐政以两江总督兼任后，盐运使是实际上管理两淮盐具体事务的盐官18。两淮盐运使的衙署名两淮都转盐运使司，简称运司，是职掌盐课税收的重要机构，其「运库存贮帑银，动盈数百万」19，时称「两淮为盐务繁区，运司为钱粮总汇，责任綦重」20。运司衙署下又分设以下四种官。

经历，一人，从七品。掌公文往来及某些具体盐务，如在北桥称放淮南各纲、食引盐，查验引数后发放给其桥旗、限票；每月定期在署堂监督收兑盐课钱粮，以及轮流委掣扬州附近之江都、甘泉二县引盐等。经历有其办公分署，名经历司，下设书吏名攒典，一人，协助经历办理事务。

知事，一人，从八品。掌淮南、淮北纲、食引盐之皮票、清单及所交银之核对、造册等。兼管江都、甘泉二县引盐之解捆，督察其缉私事务。

广盈库大使，一人，旧为「未入流」，雍正六年改为正八品。掌给发淮南、淮北纲、食盐引，定期收兑商课，经管库藏。下设攒典一人佐理。

白塔河巡检，一人，从九品。掌淮南纲、食引盐之「塔报」，即检查盐船之商名、引数，具结造报，以送有关各查验机构，以攒典一人佐理。旧设衙署白塔河巡检司，驻白塔河，后移驻盐运使衙门。与经历定期监督收兑盐课钱粮，轮流委掣

\vfill
\begin{BookNotes}
\BookNote{16.}{《清宫扬州御档选编》第5册第386-388页，王鼎、宝兴奏折。}
\BookNote{17.}{盐运使及以下诸官之设置、职掌，均据光绪重修《两淮盐法志》卷129、130〈职官门·官制上、下，不另注。}
\BookNote{18.}{佐伯富，〈清代盐政之研究（续）〉，顾南、顾学稼译文，《盐业史研究》1993年第3期，第10页。}
\BookNote{19.}{《清宫扬州御档选编》第4册第314-315页，铁保、额勒布奏折。}
\BookNote{20.}{《清宫扬州御档选编》第4册第312-313页，庐凤道德庆奏折。}
\end{BookNotes}

% 原书第 112 页（PDF 第 125 页）
\BookPage{112}{125}
\noindent 江、甘二县引盐。

以上经历、知事、库大使、巡检，同为两淮盐运使司的四首领官。

驻扬州的两淮盐政、盐运使，为管理两淮盐务的两个最大的专职盐务官。其下属或与其相关而驻扬州以外的分掌各具体盐务的官员如下。

淮南盐

淮南监掣同知，一人，正五品，驻仪征。初，仪征盐引批验所（简称「仪所」）监掣是委任府佐贰官，雍正十二年设监掣专官。乾隆元年定，所设监掣通判为专繁要缺。通判为正六品官，乾隆二十五年改设为正五品的监掣同知，遂为定制。

监掣同知掌验掣淮南引盐之政令，职司「淮南百万引盐捆、掣最要之区」，凡稽查商盐捆装、称掣，及盐船之查仓编号、验烙、催开等事，均为其重要职任21。以书吏「典吏」二人佐之。

淮南批验所大使，一人，正八品，驻仪征。专管协助监掣同知验掣淮南纲、食引盐。

淮南泰坝监掣官，一人，驻泰州。掌淮南引盐之称掣，并稽察透露夹带等弊。由两江总督会同盐政于通省佐贰官内拣选委用。以典吏二人佐理。

淮北盐

淮北监掣同知，一人，正五品，驻淮安。掌验掣淮北引盐，及淮北巡察缉私等事。以典吏一人佐之。初，淮北引盐监掣以淮安分司管理，雍正十一年改为派委府佐贰官办理。乾隆二十五年，改与淮南同，设监掣同知专官。

淮北批验所大使，一人，正八品，驻淮安。掌协助验掣淮北引盐。

\vfill
\begin{BookNotes}
\BookNote{21.}{《清宫扬州御档选编》第5册第399-400页，两江总督陶澍、江苏巡抚林则徐奏折。}
\end{BookNotes}

% 原书第 113 页（PDF 第 126 页）
\BookPage{113}{126}
淮北乌沙河巡检司巡检，一人，从九品，初驻安东坝，称安东巡检，乾隆二十九年移驻乌沙河。后移驻淮安。掌察核盐船盐引数目及督催行船等事。设书吏攒典一人佐之。

淮北永丰坝监掣官，一人，职掌同淮南泰坝监掣官。非实缺，为委任，驻清河。

盐运司分司官

两淮盐下分三个产盐区，每一产盐区设一个分司，这三个分司为:泰州分司、通州分司、海州分司，每一分司设运判一人，从六品，以典吏二人佐理。泰州分司运判驻东台，通州分司运判驻石港，海州分司旧设淮安，称淮安分司。乾隆二十四年改驻板浦，二十八年改称海州分司。

每分司下辖若干盐场，每盐场设盐课司大使一人，旧为「未入流」，雍正六年升为正八品，职掌也扩大。盐课大使各以攒典一人佐理。盐课大使各掌本盐场之制盐、交易、征课、缉私等事²²，此外，还受理盐场一般词讼，负责管理盐场水利，维护盐场地方社会治安，赈济灾荒及教育等事23。

各分司运判则总理下辖诸盐场之事，以佐理其上司盐运使。

盐务从管理灶户制盐，到盐商购盐，盐引与实际斤数的查验，及行销，中经多道环节盐务官的管理、稽核，还须各省督抚及盐法道的地方性配合，异常复杂。以行销最广的淮南盐为例，其盐船须先经由泰州坝监掣官、白塔河巡检司的两道检验，再至北桥逐船抽引掣验，而后发给其桥旗、限票，复至仪征盐引批验所，将引数、实际斤数核实，办完一切结算手续，方可售卖。

巡盐御史，清初主要任用汉人，间用旗人。大致康熙中期以后，旗人所占比例

\vfill
\begin{BookNotes}
\BookNote{22.}{光绪重修《两淮盐法志》卷130〈职官门·官制下〉。并见乾隆《历代职官表》卷61〈盐政〉，下册第1174页下，上海古籍出版社影印本。}
\BookNote{23.}{见前引何峰〈明清淮南盐区盐场大使的设置、职责及其与州县官的关系〉一文所述。}
\end{BookNotes}

% 原书第 114 页（PDF 第 127 页）
\BookPage{114}{127}
\noindent 较大，尤其是内务府包衣旗人。两淮盐政也是如此，汉人所任很少24。盐政由皇帝「特旨简充」25，康熙七年，曾实行满汉复职，满人逐渐增多，且规定带翻译满汉文的旗人笔帖式一人协助，康熙十年虽取消满汉复职，只差一人任巡盐御史，满汉参用，但笔帖式仍照旧配置，直至雍正元年取滑。巡盐御史虽按制规定一年更代，但由于是皇帝派任，所以又有连任二三年，或隔一二年又兼任的，曹雪芹的祖父曹寅，以及其姻亲同属皇室包衣奴才的李煦，都是任过四五次的26，这明显是皇帝以肥缺私惠其皇室奴才。

道光十年，裁两淮盐政，由两江总督兼理该职。次年，淮北实行票盐法，废引盐法，三十年，淮南也实行票盐法，其制度变化，此处从略。

\BookSubheading{二、扬州盐务官的其他职掌与活动}

扬州的盐务官不仅管理盐务，而且承担其他方面事务。如管理行宫，天宁、高旻、金山、焦山四处行宫，两淮盐政除咨两江总督派兵巡护外，还负责派人修葺、洒扫、草木浇灌、定时上香等事27；再如，与两江总督筹计皇帝南巡事宜28；建造寺院29；奉旨造办玉器、刊刻典籍等等30；乾隆年间修四库全书，在扬州之征集图书，以及修成后，收藏于扬州大观堂之文汇阁、镇江金山寺之文宗阁之事，也是乾隆皇帝先后交与两淮盐政李质颖、伊龄阿办理的31。

盐业促成了扬州城市的繁华及文化的发达，其文化的发达与繁荣，又凸出体现在书画、戏曲、文艺及诗文、经籍的编纂、刊刻，与文人墨客的交往等方面，而

\vfill
\begin{BookNotes}
\BookNote{24.}{前引吴静芳〈清代前期巡视两淮盐政官员之探讨〉一文有系统的统计，请参考。}
\BookNote{25.}{《清朝通典》卷35〈职官·盐政〉，第2213页，商务印书馆十通本。}
\BookNote{26.}{见光绪重修《两淮盐法志》卷131〈职官门·职名表·国朝〉的「巡盐御史」。}
\BookNote{27.}{《清宫扬州御档选编》第2册第111-112页，两淮盐政高恒奏折。}
\BookNote{28.}{《清宫扬州御档选编》第2册第91页，署理两淮盐政吉庆奏折。}
\BookNote{29.}{《清宫扬州御档选编》第3册第209页，乾隆四十八年五月二十八日上谕。}
\BookNote{30.}{《清宫扬州御档选编》第2册第146页，两淮盐政李质颖奏折。第4册第327-330页，两淮盐政阿克当阿奏折。}
\BookNote{31.}{《清宫扬州御档选编》第3册第169-172页，李质颖奏折；第207-208页，乾隆四十七年七月初八日上谕。}
\end{BookNotes}

% 原书第 115 页（PDF 第 128 页）
\BookPage{115}{128}
\noindent 这种文化特色，也影响着盐务官，不少盐务官也具有这方面的文化特色。如乾隆后期的盐政满人伊龄阿「工诗，画法梅花道人，巡盐两淮时，于扇面画梅花、兰、竹，称逸品」。扬州戏曲的发展与驻扬州盐务官有很大关系，《扬州画舫录》载：「两淮盐务例蓄花、雅两部，以备大戏。」有不少戏班，盐务官还「自制戏具，谓之内班行头」。盐政伊龄阿还「奉旨于扬州设局改曲剧」，共四年完成，名学者凌廷堪等多人以及盐务官分司运判、经历、场大使等，均参预其事32。卢见曾尤为著名，他于乾隆前期在扬州任盐运使时，「四方名流咸集，极一时文酒之盛。金农、陈撰、厉鹗、惠栋、沈大成、陈章等前后数十人，皆为上客。尝校刊《干凿度》、高氏《战国策》、郑氏《尚书大传》、李鼎祚《周易集解》及子、史等书。又补刊朱彝尊《经义考》…百余年来所罕见」33。戴震来扬州，卢见曾延其为座上客。扬州八怪中除上述金农外，卢见曾与郑板桥及高凤翰也来往密切，「唱和甚多」。还有不少工诗文、书画之名人被卢见曾延请到盐运使衙署从事文化活动34。名学者汪中，则被盐政全德聘请而「经理金山御书楼，为经史之学」。程兆熊，工诗词、书画，受盐政高恒之请，「为之写固哉亭集」35。有些盐务官还在扬州兴办书院，如安定书院就是康熙元年「巡盐御史胡文学创始」，雍正年间，盐运使尹会一又「增置学舍，为郡士肄业之所」。梅花书院为刘重选所建，其后归之有司，与安定书院「均归盐务（官）延师掌院」36。因此可以说，扬州的盐务官对扬州文化事业的发展起了一定的作用。王万祥〈清代两淮盐官文学活动研究〉一文37，详细介绍了盐官在扬州期间的文学活动、对扬州文化事业的贡献，高度评价了两淮盐务官在这方面所起到的重要作用。

另一方面，扬州的某些盐务官还担任皇帝所交与的特殊政治性事务。如刺探政情，康熙后期，两淮巡盐御史李煦曾向皇帝奏报浙江海盗之事，他「一面星飞差人到该地方确查外，谨先据风闻入告，容臣煦访确之日，再具折奏闻」。朱批：

\vfill
\begin{BookNotes}
\BookNote{32.}{李斗，〈扬州画舫录〉，第42、103、107、109页，江苏广陵古籍刻印社，1984年版。}
\BookNote{33.}{《清史列传》卷71〈文苑传二·卢见曾〉，第18册5838-5839页，中华书局，1987年版。}
\BookNote{34.}{《扬州画舫录》第218-225页。}
\BookNote{35.}{《扬州画舫录》第145、267页。}
\BookNote{36.}{《扬州画舫录》第60、62页。}
\BookNote{37.}{该文为扬州大学硕士学位论文，2010年5月。}
\end{BookNotes}

% 原书第 116 页（PDF 第 129 页）
\BookPage{116}{129}
\noindent 「知道了，再访。」38康熙五十一年，穆和伦、张廷枢审理科场案，李煦又报告扬州人对案件处理结果的反映39。扬州在清初满族征服中原时，曾发生过惨烈的「扬州十日」事件，扬州也因而成为满族皇帝关注的敏感之地，该地繁盛的戏曲对官绅百姓影响甚大，其中关涉满汉关系、南宋与金关系的内容，就曾引起乾隆皇帝的注视，认为此等内容「殊有关系」，先后命两淮盐政伊龄阿、图明阿及全德检查并删改戏曲中的这方面内容，并指示他们办理此事要「不动声色，妥协办理，不得过当滋扰」40。先是伊龄阿将《精忠传》等五种删改、抽掣，进呈乾隆帝。继之，图明阿又设局收罗「官、商及戏班教习等缴到戏曲…计二百八十四种」，将其「逐一查检」，遵照「上谕，悉心体会，窃以凡系明季国初之事，有关涉本朝字句，及南宋与金朝剧本扮演失实者，皆当遵旨酌拟删改、抽掣，另缮清本，仍将粘签之原本一并进呈，恭候钦定」41。乾隆后期，清高宗为宣扬满汉一体的忠君为国思想，改变对明清鼎革时期人物的评价，贬抑贰臣，肯定某些忠于南明而抗清的人物，其中最重要的，就是保卫扬州不屈而死的史可法。乾隆帝为了向扬州官民也是向全国官民表达他的政治观念，褒扬史可法，增加人们对满族皇帝的好感，特为史可法画像、札稿合卷（包括史可法覆摄政王书）作御制诗章，并命廷臣和韵，装璜匣贮，送往扬州，令盐政寅着放置于扬州梅花岭史可法祠中，并命其将卷内所有诗文、画像、札稿做石刻以垂永久，书卷听人展阅，但当加意保护，勿致污损42。以上这些职掌及活动，也是认识与评价两淮盐务官所应一并考虑的。

\vfill
\begin{BookNotes}
\BookNote{38.}{《清宫扬州御档选编》第1册第35-36页，李煦奏折。}
\BookNote{39.}{《清宫扬州御档选编》第1册第42页，李煦奏折。}
\BookNote{40.}{《清宫扬州御档选编》第3册第202页，图明阿奏折。}
\BookNote{41.}{《清宫扬州御档选编》第3册第198-200页，图明阿奏折，并见第200-203页，图明阿两折。}
\BookNote{42.}{《清宫扬州御档选编》第3册第185页，乾隆四十二年六月二十一日上谕。}
\end{BookNotes}
